%%%
%%% Radical "⾧"
%%%
\section*{Radical 168: ``⾧'' (长、镸)}\addcontentsline{toc}{section}{Radical 168: ⾧,长、镸}\addcontentsline{loh}{figure}{\#\#\#\# 168: ⾧}

%%%%%%%%%% 长 %%%%%%%%%%
\subsection*{长}\addcontentsline{loh}{figure}{长}

\begin{Entry}{长}{4}{⾧}[Kangxi 168]
  \begin{Phonetics}{长}{chang2}[][HSK 2]
    \definition*{s.}{Sobrenome: Chang}
    \definition{adj.}{longo; comprido; a distância entre uma extremidade e outra é grande | para sempre; duradouro; a distância entre o ponto inicial e o ponto final de um determinado período é grande | excesso; o que sobra; o que é desnecessário}
    \definition{s.}{comprimento; na direção horizontal, a distância percorrida por um objeto de uma extremidade à outra | ponto forte; especialidade; especialização em determinada área; vantagem}
    \definition{v.}{ser bom em; ser forte em; ser muito bom em algo; ser especialista em algo}
  \seealsoref{久}{jiu3}
  \synonymref{精}{jing1}
  \synonymref{久}{jiu3}
  \synonymref{善}{shan4}
  \synonymref{剩}{sheng4}
  \synonymref{优}{you1}
  \synonymref{余}{yu2}
  \synonymref{专}{zhuan1}
  \antonymref{短}{duan3}
  \antonymref{劣}{lie4}
  \antonymref{缺}{que1}
  \antonymref{弱}{ruo4}
  \antonymref{少}{shao3}
  \antonymref{生}{sheng1}
  \antonymref{幼}{you4}
  \antonymref{拙}{zhuo1}
  \end{Phonetics}
  \begin{Phonetics}{长}{zhang3}[][HSK 2,6]
    \definition{adj.}{mais velho; ancião; sênior; mais antigo; o primeiro da fila}
    \definition{s.}{ancião; pessoas mais velhas ou de posição social mais elevada | chefe; líder; dirigente; responsável}
    \definition{v.}{crescer; desenvolver"-se | surgir; começar a crescer; formar"-se; nascer em pessoas, animais, plantas ou objetos (algo) | adquirir; melhorar; aumentar; aumentar conhecimento, capacidade, etc.; tornar"-se cada vez mais ou cada vez melhor}
  \seealsoref{久}{jiu3}
  \synonymref{久}{jiu3}
  \antonymref{短}{duan3}
  \antonymref{少}{shao4}
  \antonymref{幼}{you4}
  \end{Phonetics}
\end{Entry}

\begin{Entry}{长久}{4,3}{⾧,⼃}
  \begin{Phonetics}{长久}{chang2jiu3}[][HSK 6]
    \definition{adj.}{longo; permanente; duradouro; por muito tempo; já faz muito tempo}
  \end{Phonetics}
\end{Entry}

\begin{Entry}{长大}{4,3}{⾧,⼤}
  \begin{Phonetics}{长大}{zhang3da4}[][HSK 2]
    \definition{v.}{crescer; ser criado; um estado de maturidade física ou mental completa}
  \synonymref{成熟}{cheng2shu2}
  \synonymref{成长}{cheng2zhang3}
  \synonymref{发育}{fa1yu4}
  \synonymref{高大}{gao1da4}
  \synonymref{生长}{sheng1zhang3}
  \antonymref{矮小}{ai3xiao3}
  \antonymref{衰退}{shuai1tui4}
  \antonymref{缩小}{suo1/xiao3}
  \antonymref{萎缩}{wei3suo1}
  \end{Phonetics}
\end{Entry}

\begin{Entry}{长处}{4,5}{⾧,⼡}
  \begin{Phonetics}{长处}{chang2chu4}[][HSK 3]
    \definition[个]{s.}{forte; boas qualidades; pontos fortes; especialização em determinada área}
  \seealsoref{好处}{hao3chu5}
  \seealsoref{优点}{you1dian3}
  \synonymref{好处}{hao3chu5}
  \synonymref{强项}{qiang2xiang4}
  \synonymref{所长}{suo3zhang3}
  \synonymref{特长}{te4chang2}
  \synonymref{益处}{yi4chu4}
  \synonymref{优点}{you1dian3}
  \synonymref{优势}{you1shi4}
  \antonymref{不足}{bu4zu2}
  \antonymref{短处}{duan3chu5}
  \antonymref{劣势}{lie4shi4}
  \antonymref{缺点}{que1dian3}
  \antonymref{缺陷}{que1xian4}
  \antonymref{弱点}{ruo4dian3}
  \end{Phonetics}
\end{Entry}

\begin{Entry}{长江}{4,6}{⾧,⽔}
  \begin{Phonetics}{长江}{chang2jiang1}
    \definition*{s.}{Rio Yangtze; Rio Chang Jiang}
  \end{Phonetics}
\end{Entry}

\begin{Entry}{长达}{4,6}{⾧,⾡}
  \begin{Phonetics}{长达}{chang2 da2}[][HSK 7-9]
    \definition{v.}{estender"-se tanto quanto | alongar"-se para}
  \end{Phonetics}
\end{Entry}

\begin{Entry}{长寿}{4,7}{⾧,⼨}
  \begin{Phonetics}{长寿}{chang2shou4}[][HSK 5]
    \definition{adj.}{vida longa; longevidade}
  \end{Phonetics}
\end{Entry}

\begin{Entry}{长足}{4,7}{⾧,⾜}
  \begin{Phonetics}{长足}{chang2zu2}[][HSK 7-9]
    \definition{adj.}{Literário: aos trancos e barrancos; substancial; progredindo rapidamente}
  \end{Phonetics}
\end{Entry}

\begin{Entry}{长远}{4,7}{⾧,⾡}
  \begin{Phonetics}{长远}{chang2yuan3}[][HSK 6]
    \definition{adj.}{longo prazo; longo alcance; duradouro; há muito tempo (referindo"-se ao futuro); crônico}
  \end{Phonetics}
\end{Entry}

\begin{Entry}{长征}{4,8}{⾧,⼻}
  \begin{Phonetics}{长征}{chang2zheng1}[][HSK 7-9]
    \definition*{s.}{A Longa Marcha (1934--1935), um importante movimento estratégico do Exército Vermelho dos Trabalhadores e Camponeses da China que conseguiu alcançar a base revolucionária no norte de Shaanxi depois de atravessar onze províncias e cobrir 25.000 li ou 12.500 quilômetros}
    \definition{s.}{expedição; longa marcha; viagem de longa distância; expedição de longa distância}
  \end{Phonetics}
\end{Entry}

\begin{Entry}{长城}{4,9}{⾧,⼟}
  \begin{Phonetics}{长城}{chang2cheng2}[][HSK 3]
    \definition*[段,座,条]{s.}{A Grande Muralha da China; também é usado como metáfora para uma força poderosa e inabalável, um obstáculo intransponível, etc.}
  \synonymref{堡垒}{bao3lei3}
  \end{Phonetics}
\end{Entry}

\begin{Entry}{长度}{4,9}{⾧,⼴}
  \begin{Phonetics}{长度}{chang2du4}[][HSK 5]
    \definition{s.}{comprimento; extensão; distância entre dois pontos}
  \end{Phonetics}
\end{Entry}

\begin{Entry}{长相}{4,9}{⾧,⽬}
  \begin{Phonetics}{长相}{zhang3xiang4}[][HSK 7-9]
    \definition{s.}{aparência; características; aspecto; a aparência do rosto de uma pessoa}
  \synonymref{外表}{wai4biao3}
  \synonymref{外貌}{wai4mao4}
  \end{Phonetics}
\end{Entry}

\begin{Entry}{长效}{4,10}{⾧,⽁}
  \begin{Phonetics}{长效}{chang2xiao4}[][HSK 7-9]
    \definition{adj.}{duradouro}
    \definition{s.}{efeito duradouro}
    \definition{v.}{ser eficaz por um período prolongado}
  \end{Phonetics}
\end{Entry}

\begin{Entry}{长途}{4,10}{⾧,⾡}
  \begin{Phonetics}{长途}{chang2tu2}[][HSK 4]
    \definition{adj.}{de longa distância; longe}
    \definition[段,次,程]{s.}{longa"-distância; referindo"-se especificamente a chamadas telefônicas de longa distância ou ônibus de longa distância}
  \end{Phonetics}
\end{Entry}

\begin{Entry}{长假}{4,11}{⾧,⼈}
  \begin{Phonetics}{长假}{chang2jia4}[][HSK 6]
    \definition{s.}{longa licença de ausência | feriado prolongado | Obsoleto: renúncia | férias longas | refere"-se a um feriado nacional de uma semana na China, começando em 1º de maio e 1º de outubro}
  \end{Phonetics}
\end{Entry}

\begin{Entry}{长颈鹿}{4,11,11}{⾧,⾴,⿅}
  \begin{Phonetics}{长颈鹿}{chang2jing3lu4}
    \definition[只,头,群]{s.}{girafa (mamífero)}
  \end{Phonetics}
\end{Entry}

\begin{Entry}{长期}{4,12}{⾧,⽉}
  \begin{Phonetics}{长期}{chang2qi1}[][HSK 3]
    \definition{adj.}{secular; longo prazo; longo alcance; durante um longo período de tempo}
    \definition{s.}{longo prazo; por muito tempo}
  \synonymref{长久}{chang2jiu3}
  \synonymref{长远}{chang2yuan3}
  \synonymref{持久}{chi2jiu3}
  \synonymref{永久}{yong3jiu3}
  \synonymref{悠久}{you1jiu3}
  \antonymref{短期}{duan3qi1}
  \antonymref{临时}{lin2shi2}
  \antonymref{瞬间}{shun4jian1}
  \antonymref{暂时}{zan4shi2}
  \end{Phonetics}
\end{Entry}

\begin{Entry}{长期以来}{4,12,4,7}{⾧,⽉,⼈,⽊}
  \begin{Phonetics}{长期以来}{chang2qi1yi3lai2}[][HSK 7-9]
    \definition{adv.}{por muito tempo passado | desde muito tempo atrás}
  \end{Phonetics}
\end{Entry}

\begin{Entry}{长短}{4,12}{⾧,⽮}
  \begin{Phonetics}{长短}{chang2duan3}[][HSK 6]
    \definition{s.}{comprimento | acidente; infortúnio | certo e errado; pontos fortes e fracos}
  \end{Phonetics}
\end{Entry}

\begin{Entry}{长跑}{4,12}{⾧,⾜}
  \begin{Phonetics}{长跑}{chang2pao3}[][HSK 6]
    \definition{s.}{corrida de longa distância}
  \antonymref{短跑}{duan3pao3}
  \end{Phonetics}
\end{Entry}

\begin{Entry}{长辈}{4,12}{⾧,⾞}
  \begin{Phonetics}{长辈}{zhang3bei4}[][HSK 7-9]
    \definition[位,名,个,些]{s.}{ancião; sênior; membro mais velho de uma família}
  \synonymref{前辈}{qian2bei4}
  \end{Phonetics}
\end{Entry}

%%%%% EOF %%%%%

